%% kapa: exam=2 round=32 year=2021 period=2 subject=appraisal_theory subject_ko=감정평가이론 seq=3 number=3 points=20
%% source: Q-Net 감정평가사 기출문제 (cst00309) artlSeq=5208993 file=raw/1008/5208993/감평 2차 시험지(2021).pdf
%% conversion: kordoc 4.24.51 -> tools/gen2.py (auto); review=proofread-sonnet; flags=none
\begin{longquestion}{【문제 3】}
광평수(廣坪數) 토지란 해당 토지가 속해 있는 시장지역에서 일반적으로 사용하는 표준적 규모보다 훨씬 더 크다고 인식되는 토지로서, 최근에 대단위 아파트 단지개발 및 복합용도개발 등으로 인해 광평수 토지에 대한 감정평가가 증가하고 있다. 이와 관련한 다음 물음에 답하시오. (20점) \par
\medskip\noindent 물음 1) 광평수 토지면적이 해당 토지의 가치에 미치는 영향을 감가(減價)와 증가(增價)로 나누어 설명하시오. (10점) \par
\medskip\noindent 물음 2) 광평수 토지의 최유효이용이 단독이용(single use)인 경우 감정평가방법에 대해 설명하시오. (10점) \par
\end{longquestion}
